\documentclass[10pt]{article}
\usepackage[letterpaper,margin=0.62in]{geometry}
\usepackage[T1]{fontenc}
\usepackage{lmodern}
\usepackage{amsmath}
\usepackage{booktabs}
\usepackage{graphicx}
\usepackage{longtable}
\usepackage{pdflscape}
\usepackage[table]{xcolor}
\usepackage{microtype}
\usepackage{enumitem}
\usepackage[hidelinks]{hyperref}
\usepackage{fancyhdr}

\definecolor{pairwise}{HTML}{B42318}
\definecolor{unary}{HTML}{374151}
\definecolor{original}{HTML}{355F8A}
\definecolor{ink}{HTML}{111827}
\definecolor{muted}{HTML}{4B5563}
\definecolor{panel}{HTML}{F3F4F6}

\newcommand{\StepNineMinGain}{5.32\%}
\newcommand{\StepNineMaxGain}{9.60\%}
\newcommand{\StepNinePrimaryTable}{%
\begin{tabular}{lrrrr}
\toprule Dataset & $\Delta$ tok/s & 95\% CI & Relative & I/T/H \\
\midrule
GSM8K & 13.47 & [12.11, 14.81] & 8.44\% & 121/0/7 \\
MATH500 & 15.63 & [13.88, 17.42] & 9.60\% & 120/0/8 \\
HumanEval & 10.66 & [9.08, 12.28] & 7.61\% & 144/0/20 \\
MT-Bench-controlled & 7.07 & [5.53, 8.62] & 5.41\% & 67/0/13 \\
MT-Bench-native & 7.20 & [5.43, 9.02] & 5.32\% & 68/0/12 \\
\bottomrule
\end{tabular}}


\pagestyle{fancy}
\fancyhf{}
\lhead{\footnotesize Step 9: Qwen3-4B throughput validation}
\rhead{\footnotesize Frozen benchmark report}
\cfoot{\footnotesize\thepage}
\setlength{\headheight}{14pt}
\setlength{\parindent}{0pt}
\setlength{\parskip}{4pt}
\setlist[itemize]{leftmargin=16pt,itemsep=2pt,topsep=2pt}
\renewcommand{\arraystretch}{1.08}
\newcommand{\pairwise}{\textcolor{pairwise}{Pairwise-K16}}
\newcommand{\unaryname}{\textcolor{unary}{unary DDTree}}

\begin{document}

\begin{center}
{\LARGE\bfseries Conditional Draft Trees with DFlash2}\\[3pt]
{\large One-H100 Qwen3-4B throughput benchmark}\\[5pt]
{\small Frozen protocol commit
\texttt{1bfb1cb8cf3cb40308d520de608a01fe5a86eafe};
results commit \texttt{84c839b3df84068a5ec010d8b64550f4dbf40457}}
\end{center}

\vspace{2pt}
\colorbox{panel}{%
\parbox{\dimexpr\linewidth-2\fboxsep\relax}{%
\textbf{Result in one sentence.}
Within the same DFlash2 system, predecessor-conditioned \pairwise{} allocation
improves end-to-end throughput over released \unaryname{} by
\StepNineMinGain--\StepNineMaxGain{} across every tested slice; however,
released original DFlash and original DFlash+DDTree remain faster in absolute
tokens/s.}}

\section{Research question and experimental boundary}

The controlled question is:
\emph{Given the same DFlash2 checkpoint, top-16 candidate lattice, packed
target verifier, prompt, and verification-node budget, does
predecessor-conditioned tree allocation improve end-to-end output
throughput over released unary DDTree allocation?}

The practical system question is broader:
\emph{Does DFlash2 with Pairwise-DDTree outperform raw released DFlash and
released DFlash+DDTree in absolute tokens/s?}
The first comparison isolates allocation. The second changes the drafter
checkpoint, architecture, and draft attention backend, so it is descriptive
rather than a causal allocator ablation.

\begin{center}
\StepNinePrimaryTable
\end{center}

All five primary paired prompt/conversation-bootstrap 95\% confidence
intervals exclude zero. The gain is positive on 520 of 580 prompt clusters.

\begin{figure}[ht]
\centering
\includegraphics[width=0.84\linewidth]{figures/step9_primary_gain.pdf}
\caption{Primary controlled comparison at \(B=16\). Error bars are paired
prompt/conversation-bootstrap 95\% confidence intervals over end-to-end
output tokens/s, including target prefill.}
\label{fig:primary}
\end{figure}

\section{Frozen benchmark configuration}

\begin{tabular}{@{}ll@{}}
\toprule
Component & Frozen value \\
\midrule
Target & \texttt{Qwen/Qwen3-4B} at revision \texttt{1cfa9a72...} \\
Original drafter & \texttt{z-lab/Qwen3-4B-DFlash-b16} at revision \texttt{b74e3a32...} \\
DFlash2 drafter & \texttt{mgoin/Qwen3-4B-speculator.dflash2} at revision \texttt{e3e7a18e...} \\
Hardware & One NVIDIA H100 80GB HBM3 \\
Software & Python 3.10.12; PyTorch 2.7.0; CUDA 12.8; Transformers 5.16.1 \\
Precision / decoding & BF16; greedy; maximum 256 new tokens \\
Attention & Target SDPA; original draft FlashAttention2; DFlash2 draft SDPA \\
Timing & Three repetitions per method and prompt/turn; full method-order shuffle \\
Primary endpoint & End-to-end output tokens/s, including prefill \\
Uncertainty & 10,000 paired prompt/conversation bootstrap resamples; seed 20260906 \\
\bottomrule
\end{tabular}

\vspace{5pt}
Datasets use pinned revisions and deterministic seed-0 selection: 128 GSM8K,
128 MATH500, all 164 HumanEval, and 80 two-turn MT-Bench conversations.
MT-Bench is measured with both controlled sequential-baseline history and
native per-method history. The full provenance, dataset revisions, runtime
versions, artifact hashes, and timing-source counts are recorded in
\texttt{analysis/step9-4b-throughput/provenance.json}.

\clearpage
\section{Absolute throughput: allocator win, not system win}

\begin{figure}[ht]
\centering
\includegraphics[width=\linewidth]{figures/step9_system_throughput.pdf}
\caption{Absolute end-to-end output throughput at \(B=16\). The original
DFlash family is faster than the tested DFlash2 family. Within DFlash2,
Pairwise consistently improves over unary DDTree.}
\label{fig:absolute}
\end{figure}

At \(B=16\), \pairwise{} raises DFlash2 throughput from 163.32 to
176.79 tok/s on GSM8K, 168.58 to 184.21 on MATH500, 145.77 to 156.44 on
HumanEval, 129.57 to 136.64 on controlled MT-Bench, and 130.46 to 137.66 on
native MT-Bench. The associated relative gains range from
\StepNineMinGain{} to \StepNineMaxGain.

The stronger practical baseline remains released original DFlash:
raw original DFlash reaches 152.98--268.28 tok/s, while original
DFlash+DDTree-B16 reaches 177.95--285.39 tok/s across the same slices.
Every Pairwise-B16 cross-drafter confidence interval remains below zero,
including after normalizing each family by its own sequential-target process.

\textbf{Interpretation.}
The experiment answers the narrow scientific question positively:
conditioning allocates the DFlash2 tree more efficiently and converts that
acceptance gain into throughput. It answers the product-level question
negatively for these checkpoints: the tested Pairwise-DFlash2 system is not
the fastest end-to-end system.

\section{Behavior across verification budgets}

\begin{figure}[ht]
\centering
\includegraphics[width=\linewidth]{figures/step9_budget_throughput.pdf}
\caption{End-to-end throughput across verification budgets. Fixed Unary-K16
is shown at \(B=32,64\) to separate candidate breadth from allocation.}
\label{fig:budget-throughput}
\end{figure}

\pairwise{} exceeds released DFlash2 unary DDTree at every tested
\(B\in\{7,16,32,64\}\) and on every dataset. At \(B=32,64\), released unary
DDTree receives candidate width \(K=\min(B,|V|)\), while Pairwise remains on
the native top-16 selector lattice. Its advantage therefore does not come
from wider support.

\clearpage
\section{Mechanism: acceptance and target-call efficiency}

\begin{figure}[ht]
\centering
\includegraphics[width=\linewidth]{figures/step9_budget_acceptance.pdf}
\caption{Mean matched draft tokens per verification round. Pairwise
conditioning consistently produces more useful speculative paths under the
same node budget.}
\label{fig:acceptance}
\end{figure}

At primary \(B=16\), mean matched draft tokens per round increase:
\begin{itemize}
\item GSM8K: 4.347 \(\rightarrow\) 4.823;
\item MATH500: 4.375 \(\rightarrow\) 4.904;
\item HumanEval: 3.819 \(\rightarrow\) 4.197;
\item controlled MT-Bench: 3.116 \(\rightarrow\) 3.360;
\item native MT-Bench: 3.109 \(\rightarrow\) 3.360.
\end{itemize}

More matched and committed draft tokens reduce target calls. Pairwise does
add conditional tree-construction work, but the complete timing decomposition
in Appendix~\ref{tab:all-timing} shows that the saved target work dominates
under the packed Qwen3-4B verifier.

\clearpage
\section{Correctness and task quality}

\begin{figure}[ht]
\centering
\begin{minipage}[t]{0.49\linewidth}
\centering
\includegraphics[width=\linewidth]{figures/step9_exact_output.pdf}
\caption{Full-output exact match against one-token sequential BF16 decoding.}
\label{fig:exact}
\end{minipage}
\hfill
\begin{minipage}[t]{0.49\linewidth}
\centering
\includegraphics[width=\linewidth]{figures/step9_task_quality.pdf}
\caption{Direct task quality at \(B=16\). Differences are descriptive.}
\label{fig:quality}
\end{minipage}
\end{figure}

BF16 packed/block verification is not token-exact with one-token sequential
target decoding on every prompt. At \(B=16\), Pairwise exact-output rates are
57.0\% on GSM8K, 35.2\% on MATH500, 45.1\% on HumanEval, 36.3\% on controlled
MT-Bench, and 37.5\% on native MT-Bench. Many outputs hit the 256-token cap,
but the result still does not establish strict lossless equivalence.

Direct task scores do not show a consistent Pairwise-B16 regression:
GSM8K is 62/128 for Pairwise versus 64/128 for unary, MATH500 is 18/128
for both, and HumanEval is 89/164 versus 88/164. These are not powered
equivalence tests. MT-Bench responses were exported but not externally
judged, so no MT-Bench quality score is claimed.

\clearpage
\section{Conclusions and limitations}

\textbf{Supported conclusions}
\begin{itemize}
\item Under the tested one-H100 Qwen3-4B setup, predecessor-conditioned
Pairwise allocation improves both matched draft tokens and end-to-end
throughput over released unary DDTree allocation on the same DFlash2
checkpoint.
\item The controlled effect holds across math, code, and dialogue tasks,
both MT-Bench trajectory regimes, and every tested node budget.
\item Wider unary support does not explain the gain: Pairwise-K16 remains
ahead when released unary DDTree receives \(K=32\) or \(K=64\).
\end{itemize}

\textbf{Important boundaries}
\begin{itemize}
\item Released original DFlash and original DFlash+DDTree are still faster
in absolute tokens/s. Pairwise improves DFlash2; it does not win the complete
cross-drafter system comparison.
\item Cross-drafter rows conflate model architecture, checkpoint quality,
implementation, and FlashAttention2-versus-SDPA draft attention.
\item BF16 numerical trajectories are not universally exact, and the task
sample sizes do not prove quality equivalence.
\item Peak reserved memory includes allocator caching; peak allocated memory
is the more meaningful method-level footprint.
\item This is one GPU type, one target checkpoint, one maximum output length,
and greedy decoding. Generalization beyond this setup requires measurement.
\end{itemize}

\section{Reproduction and provenance}

Figures and appendix tables are regenerated only from frozen local analysis:
\begin{verbatim}
python research_notes/figures/generate_step9_benchmark_report.py
cd research_notes
pdflatex -interaction=nonstopmode -halt-on-error \
  -output-directory=build_step9_report step9_4b_benchmark_report.tex
pdflatex -interaction=nonstopmode -halt-on-error \
  -output-directory=build_step9_report step9_4b_benchmark_report.tex
\end{verbatim}

The PDF is generated from CSV/JSON analysis, not manually transcribed chart
values. Raw benchmark artifacts are retained locally under
\texttt{artifacts/step9\_4b/1bfb1.../}; compact evidence and hashes are
committed under \texttt{analysis/step9-4b-throughput/}. All 138 remote
benchmark, analysis, quality, and log files were copied locally and verified
against remote SHA-256 hashes before the H100 was released.

\appendix
\clearpage
\section{Complete frozen result tables}

The following tables reproduce every row in the frozen method metrics,
paired-comparison, timing-decomposition, correctness, and directly scored
quality outputs. NaN/not-applicable values are rendered as dashes.

\begin{landscape}
\scriptsize
\setlength{\tabcolsep}{3pt}

\begin{longtable}{lllrrrrrrr}
\caption{Complete per-method throughput and resource metrics.}\label{tab:all-throughput}\\
\toprule
Dataset & Family & Method & B & E2E tok/s & Decode tok/s & Speedup & Target calls & TTFT ms & Peak GiB \\
\midrule
\endfirsthead
\toprule
Dataset & Family & Method & B & E2E tok/s & Decode tok/s & Speedup & Target calls & TTFT ms & Peak GiB \\
\midrule
\endhead
\midrule
\multicolumn{10}{r}{Continued on next page}\\
\endfoot
\bottomrule
\endlastfoot
GSM8K & dflash2 & Sequential-Baseline & -- & 49.70 & 49.92 & -- & 237.6 & 20.32 & 10.27 \\
GSM8K & dflash2 & Raw-DFlash2 & -- & 155.74 & 158.79 & 3.13 & 50.2 & 20.55 & 10.29 \\
GSM8K & dflash2 & DFlash2-Original-DDTree-B16 & 16 & 163.32 & 166.76 & 3.29 & 44.8 & 20.56 & 10.30 \\
GSM8K & dflash2 & DFlash2-Original-DDTree-B32 & 32 & 169.27 & 172.96 & 3.41 & 42.3 & 20.46 & 10.32 \\
GSM8K & dflash2 & DFlash2-Original-DDTree-B64 & 64 & 164.14 & 167.61 & 3.30 & 40.9 & 20.47 & 10.35 \\
GSM8K & dflash2 & DFlash2-Original-DDTree-B7 & 7 & 151.98 & 154.96 & 3.06 & 49.3 & 20.52 & 10.29 \\
GSM8K & dflash2 & DFlash2-Pairwise-K16-B16 & 16 & 176.79 & 180.84 & 3.56 & 41.1 & 20.69 & 10.30 \\
GSM8K & dflash2 & DFlash2-Pairwise-K16-B32 & 32 & 185.24 & 189.66 & 3.73 & 38.7 & 20.57 & 10.31 \\
GSM8K & dflash2 & DFlash2-Pairwise-K16-B64 & 64 & 189.02 & 193.64 & 3.80 & 37.3 & 20.50 & 10.35 \\
GSM8K & dflash2 & DFlash2-Pairwise-K16-B7 & 7 & 158.35 & 161.57 & 3.19 & 47.1 & 20.48 & 10.29 \\
GSM8K & dflash2 & DFlash2-Unary-K16-B32 & 32 & 171.14 & 174.92 & 3.44 & 42.3 & 20.55 & 10.32 \\
GSM8K & dflash2 & DFlash2-Unary-K16-B64 & 64 & 173.76 & 177.65 & 3.50 & 41.0 & 20.49 & 10.35 \\
GSM8K & original & Sequential-Baseline & -- & 53.07 & 53.31 & -- & 237.6 & 19.48 & 8.64 \\
GSM8K & original & DFlash-Original-DDTree-B16 & 16 & 261.22 & 268.68 & 4.92 & 34.0 & 19.54 & 8.68 \\
GSM8K & original & DFlash-Original-DDTree-B32 & 32 & 278.95 & 287.44 & 5.26 & 31.6 & 19.53 & 8.69 \\
GSM8K & original & DFlash-Original-DDTree-B64 & 64 & 294.45 & 303.96 & 5.55 & 29.9 & 19.46 & 8.72 \\
GSM8K & original & Raw-DFlash & -- & 232.97 & 238.90 & 4.39 & 39.2 & 19.40 & 8.67 \\
MATH500 & dflash2 & Sequential-Baseline & -- & 50.98 & 51.19 & -- & 250.3 & 19.76 & 10.32 \\
MATH500 & dflash2 & Raw-DFlash2 & -- & 162.24 & 165.33 & 3.18 & 52.7 & 20.07 & 10.36 \\
MATH500 & dflash2 & DFlash2-Original-DDTree-B16 & 16 & 168.58 & 171.96 & 3.31 & 47.3 & 20.05 & 10.36 \\
MATH500 & dflash2 & DFlash2-Original-DDTree-B32 & 32 & 172.63 & 176.17 & 3.39 & 45.1 & 20.05 & 10.37 \\
MATH500 & dflash2 & DFlash2-Original-DDTree-B64 & 64 & 167.90 & 171.23 & 3.29 & 43.5 & 19.82 & 10.40 \\
MATH500 & dflash2 & DFlash2-Original-DDTree-B7 & 7 & 156.12 & 159.02 & 3.06 & 52.3 & 19.97 & 10.36 \\
MATH500 & dflash2 & DFlash2-Pairwise-K16-B16 & 16 & 184.21 & 188.26 & 3.61 & 43.0 & 20.07 & 10.36 \\
MATH500 & dflash2 & DFlash2-Pairwise-K16-B32 & 32 & 191.93 & 196.32 & 3.76 & 40.7 & 20.14 & 10.37 \\
MATH500 & dflash2 & DFlash2-Pairwise-K16-B64 & 64 & 196.19 & 200.74 & 3.85 & 39.2 & 19.98 & 10.39 \\
MATH500 & dflash2 & DFlash2-Pairwise-K16-B7 & 7 & 165.00 & 168.24 & 3.24 & 49.4 & 19.94 & 10.36 \\
MATH500 & dflash2 & DFlash2-Unary-K16-B32 & 32 & 174.20 & 177.80 & 3.42 & 45.2 & 20.02 & 10.37 \\
MATH500 & dflash2 & DFlash2-Unary-K16-B64 & 64 & 177.67 & 181.41 & 3.48 & 43.5 & 19.96 & 10.40 \\
MATH500 & original & Sequential-Baseline & -- & 53.37 & 53.59 & -- & 250.3 & 19.29 & 8.69 \\
MATH500 & original & DFlash-Original-DDTree-B16 & 16 & 285.39 & 294.00 & 5.34 & 33.6 & 19.54 & 8.74 \\
MATH500 & original & DFlash-Original-DDTree-B32 & 32 & 307.48 & 317.62 & 5.76 & 30.8 & 19.84 & 8.74 \\
MATH500 & original & DFlash-Original-DDTree-B64 & 64 & 326.56 & 337.72 & 6.11 & 29.2 & 19.36 & 8.77 \\
MATH500 & original & Raw-DFlash & -- & 268.28 & 275.98 & 5.02 & 37.4 & 19.37 & 8.74 \\
HumanEval & dflash2 & Sequential-Baseline & -- & 49.12 & 49.31 & -- & 254.2 & 20.58 & 10.32 \\
HumanEval & dflash2 & Raw-DFlash2 & -- & 135.44 & 137.61 & 2.76 & 61.4 & 20.89 & 10.37 \\
HumanEval & dflash2 & DFlash2-Original-DDTree-B16 & 16 & 145.77 & 148.33 & 2.97 & 53.5 & 20.76 & 10.38 \\
HumanEval & dflash2 & DFlash2-Original-DDTree-B32 & 32 & 150.82 & 153.55 & 3.07 & 50.2 & 20.76 & 10.39 \\
HumanEval & dflash2 & DFlash2-Original-DDTree-B64 & 64 & 147.76 & 150.39 & 3.01 & 47.9 & 20.70 & 10.40 \\
HumanEval & dflash2 & DFlash2-Original-DDTree-B7 & 7 & 133.83 & 136.00 & 2.73 & 59.5 & 20.80 & 10.37 \\
HumanEval & dflash2 & DFlash2-Pairwise-K16-B16 & 16 & 156.44 & 159.39 & 3.19 & 49.7 & 20.82 & 10.38 \\
HumanEval & dflash2 & DFlash2-Pairwise-K16-B32 & 32 & 165.11 & 168.41 & 3.36 & 46.1 & 20.89 & 10.39 \\
HumanEval & dflash2 & DFlash2-Pairwise-K16-B64 & 64 & 171.07 & 174.59 & 3.48 & 43.8 & 20.78 & 10.40 \\
HumanEval & dflash2 & DFlash2-Pairwise-K16-B7 & 7 & 139.45 & 141.78 & 2.84 & 56.9 & 20.69 & 10.37 \\
HumanEval & dflash2 & DFlash2-Unary-K16-B32 & 32 & 152.18 & 154.96 & 3.10 & 50.3 & 20.74 & 10.39 \\
HumanEval & dflash2 & DFlash2-Unary-K16-B64 & 64 & 157.18 & 160.16 & 3.20 & 47.9 & 20.75 & 10.40 \\
HumanEval & original & Sequential-Baseline & -- & 51.91 & 52.12 & -- & 254.2 & 19.80 & 8.69 \\
HumanEval & original & DFlash-Original-DDTree-B16 & 16 & 266.87 & 274.37 & 5.14 & 35.5 & 20.05 & 8.76 \\
HumanEval & original & DFlash-Original-DDTree-B32 & 32 & 288.37 & 297.15 & 5.55 & 32.6 & 20.06 & 8.76 \\
HumanEval & original & DFlash-Original-DDTree-B64 & 64 & 302.68 & 312.22 & 5.83 & 30.9 & 19.86 & 8.78 \\
HumanEval & original & Raw-DFlash & -- & 240.60 & 246.78 & 4.63 & 40.8 & 19.97 & 8.75 \\
MT-Bench-controlled & dflash2 & Sequential-Baseline & -- & 51.53 & 51.89 & -- & 421.9 & 19.57 & 10.37 \\
MT-Bench-controlled & dflash2 & Raw-DFlash2 & -- & 117.81 & 120.86 & 2.29 & 134.4 & 19.79 & 10.47 \\
MT-Bench-controlled & dflash2 & DFlash2-Original-DDTree-B16 & 16 & 129.57 & 133.30 & 2.52 & 111.5 & 19.78 & 10.48 \\
MT-Bench-controlled & dflash2 & DFlash2-Original-DDTree-B32 & 32 & 134.05 & 138.08 & 2.61 & 104.4 & 19.75 & 10.48 \\
MT-Bench-controlled & dflash2 & DFlash2-Original-DDTree-B64 & 64 & 130.69 & 134.60 & 2.54 & 98.7 & 19.76 & 10.49 \\
MT-Bench-controlled & dflash2 & DFlash2-Original-DDTree-B7 & 7 & 119.03 & 122.18 & 2.31 & 123.6 & 19.79 & 10.47 \\
MT-Bench-controlled & dflash2 & DFlash2-Pairwise-K16-B16 & 16 & 136.64 & 140.66 & 2.66 & 104.8 & 19.83 & 10.48 \\
MT-Bench-controlled & dflash2 & DFlash2-Pairwise-K16-B32 & 32 & 144.61 & 149.05 & 2.81 & 96.8 & 19.88 & 10.48 \\
MT-Bench-controlled & dflash2 & DFlash2-Pairwise-K16-B64 & 64 & 148.83 & 153.28 & 2.89 & 91.3 & 19.85 & 10.49 \\
MT-Bench-controlled & dflash2 & DFlash2-Pairwise-K16-B7 & 7 & 123.68 & 126.99 & 2.40 & 118.8 & 19.84 & 10.47 \\
MT-Bench-controlled & dflash2 & DFlash2-Unary-K16-B32 & 32 & 135.92 & 140.19 & 2.64 & 104.2 & 19.79 & 10.48 \\
MT-Bench-controlled & dflash2 & DFlash2-Unary-K16-B64 & 64 & 139.98 & 144.43 & 2.72 & 98.6 & 19.81 & 10.49 \\
MT-Bench-controlled & original & Sequential-Baseline & -- & 53.81 & 54.19 & -- & 421.9 & 19.24 & 8.73 \\
MT-Bench-controlled & original & DFlash-Original-DDTree-B16 & 16 & 183.57 & 190.26 & 3.41 & 99.5 & 19.24 & 8.85 \\
MT-Bench-controlled & original & DFlash-Original-DDTree-B32 & 32 & 197.18 & 204.97 & 3.66 & 92.3 & 19.45 & 8.85 \\
MT-Bench-controlled & original & DFlash-Original-DDTree-B64 & 64 & 210.06 & 219.01 & 3.90 & 86.8 & 19.30 & 8.87 \\
MT-Bench-controlled & original & Raw-DFlash & -- & 158.29 & 163.53 & 2.94 & 126.0 & 19.30 & 8.85 \\
MT-Bench-native & dflash2 & Sequential-Baseline & -- & 51.90 & 52.27 & -- & 421.9 & 19.54 & 10.37 \\
MT-Bench-native & dflash2 & Raw-DFlash2 & -- & 119.17 & 122.27 & 2.30 & 134.8 & 19.70 & 10.47 \\
MT-Bench-native & dflash2 & DFlash2-Original-DDTree-B16 & 16 & 130.46 & 134.22 & 2.52 & 111.7 & 19.76 & 10.48 \\
MT-Bench-native & dflash2 & DFlash2-Original-DDTree-B32 & 32 & 135.35 & 139.43 & 2.61 & 104.1 & 19.69 & 10.48 \\
MT-Bench-native & dflash2 & DFlash2-Original-DDTree-B64 & 64 & 132.24 & 136.19 & 2.55 & 98.0 & 19.71 & 10.49 \\
MT-Bench-native & dflash2 & DFlash2-Original-DDTree-B7 & 7 & 120.34 & 123.55 & 2.32 & 123.7 & 19.82 & 10.47 \\
MT-Bench-native & dflash2 & DFlash2-Pairwise-K16-B16 & 16 & 137.66 & 141.73 & 2.66 & 104.8 & 19.81 & 10.48 \\
MT-Bench-native & dflash2 & DFlash2-Pairwise-K16-B32 & 32 & 146.03 & 150.53 & 2.82 & 96.2 & 19.84 & 10.48 \\
MT-Bench-native & dflash2 & DFlash2-Pairwise-K16-B64 & 64 & 150.94 & 155.72 & 2.91 & 91.5 & 19.83 & 10.49 \\
MT-Bench-native & dflash2 & DFlash2-Pairwise-K16-B7 & 7 & 124.85 & 128.21 & 2.41 & 119.1 & 19.72 & 10.47 \\
MT-Bench-native & dflash2 & DFlash2-Unary-K16-B32 & 32 & 136.71 & 141.02 & 2.64 & 104.5 & 19.73 & 10.48 \\
MT-Bench-native & dflash2 & DFlash2-Unary-K16-B64 & 64 & 140.56 & 145.08 & 2.71 & 98.3 & 19.87 & 10.49 \\
MT-Bench-native & original & Sequential-Baseline & -- & 52.08 & 52.46 & -- & 421.9 & 19.93 & 8.73 \\
MT-Bench-native & original & DFlash-Original-DDTree-B16 & 16 & 177.95 & 184.46 & 3.42 & 99.1 & 19.96 & 8.85 \\
MT-Bench-native & original & DFlash-Original-DDTree-B32 & 32 & 191.02 & 198.53 & 3.67 & 92.3 & 20.14 & 8.85 \\
MT-Bench-native & original & DFlash-Original-DDTree-B64 & 64 & 204.24 & 212.91 & 3.92 & 86.2 & 20.04 & 8.87 \\
MT-Bench-native & original & Raw-DFlash & -- & 152.98 & 158.01 & 2.94 & 124.9 & 20.00 & 8.85 \\
\end{longtable}


\begin{longtable}{lllrrrrr}
\caption{Complete acceptance and exact-output metrics.}\label{tab:all-acceptance}\\
\toprule
Dataset & Family & Method & B & Matched/round & Committed/round & Exact \% & Truncated \% \\
\midrule
\endfirsthead
\toprule
Dataset & Family & Method & B & Matched/round & Committed/round & Exact \% & Truncated \% \\
\midrule
\endhead
\midrule
\multicolumn{8}{r}{Continued on next page}\\
\endfoot
\bottomrule
\endlastfoot
GSM8K & dflash2 & Sequential-Baseline & -- & 0.000 & 1.000 & -- & 62.5 \\
GSM8K & dflash2 & Raw-DFlash2 & -- & 3.805 & 4.787 & 49.2 & 64.8 \\
GSM8K & dflash2 & DFlash2-Original-DDTree-B16 & 16 & 4.347 & 5.328 & 57.8 & 63.3 \\
GSM8K & dflash2 & DFlash2-Original-DDTree-B32 & 32 & 4.660 & 5.638 & 56.2 & 62.5 \\
GSM8K & dflash2 & DFlash2-Original-DDTree-B64 & 64 & 4.844 & 5.822 & 53.1 & 61.7 \\
GSM8K & dflash2 & DFlash2-Original-DDTree-B7 & 7 & 3.871 & 4.855 & 50.8 & 62.5 \\
GSM8K & dflash2 & DFlash2-Pairwise-K16-B16 & 16 & 4.823 & 5.801 & 57.0 & 63.3 \\
GSM8K & dflash2 & DFlash2-Pairwise-K16-B32 & 32 & 5.179 & 6.155 & 50.8 & 61.7 \\
GSM8K & dflash2 & DFlash2-Pairwise-K16-B64 & 64 & 5.396 & 6.372 & 55.5 & 62.5 \\
GSM8K & dflash2 & DFlash2-Pairwise-K16-B7 & 7 & 4.095 & 5.075 & 52.3 & 62.5 \\
GSM8K & dflash2 & DFlash2-Unary-K16-B32 & 32 & 4.667 & 5.645 & 57.0 & 62.5 \\
GSM8K & dflash2 & DFlash2-Unary-K16-B64 & 64 & 4.837 & 5.816 & 52.3 & 61.7 \\
GSM8K & original & Sequential-Baseline & -- & 0.000 & 1.000 & -- & 62.5 \\
GSM8K & original & DFlash-Original-DDTree-B16 & 16 & 6.261 & 7.261 & 55.5 & 63.3 \\
GSM8K & original & DFlash-Original-DDTree-B32 & 32 & 6.821 & 7.821 & 50.8 & 64.8 \\
GSM8K & original & DFlash-Original-DDTree-B64 & 64 & 7.229 & 8.229 & 59.4 & 61.7 \\
GSM8K & original & Raw-DFlash & -- & 5.342 & 6.342 & 55.5 & 62.5 \\
MATH500 & dflash2 & Sequential-Baseline & -- & 0.000 & 1.000 & -- & 89.1 \\
MATH500 & dflash2 & Raw-DFlash2 & -- & 3.869 & 4.854 & 44.5 & 88.3 \\
MATH500 & dflash2 & DFlash2-Original-DDTree-B16 & 16 & 4.375 & 5.357 & 46.9 & 88.3 \\
MATH500 & dflash2 & DFlash2-Original-DDTree-B32 & 32 & 4.623 & 5.604 & 45.3 & 88.3 \\
MATH500 & dflash2 & DFlash2-Original-DDTree-B64 & 64 & 4.820 & 5.801 & 39.1 & 88.3 \\
MATH500 & dflash2 & DFlash2-Original-DDTree-B7 & 7 & 3.879 & 4.862 & 47.7 & 88.3 \\
MATH500 & dflash2 & DFlash2-Pairwise-K16-B16 & 16 & 4.904 & 5.883 & 35.2 & 87.5 \\
MATH500 & dflash2 & DFlash2-Pairwise-K16-B32 & 32 & 5.234 & 6.212 & 39.8 & 87.5 \\
MATH500 & dflash2 & DFlash2-Pairwise-K16-B64 & 64 & 5.456 & 6.434 & 41.4 & 88.3 \\
MATH500 & dflash2 & DFlash2-Pairwise-K16-B7 & 7 & 4.155 & 5.139 & 42.2 & 88.3 \\
MATH500 & dflash2 & DFlash2-Unary-K16-B32 & 32 & 4.619 & 5.599 & 43.8 & 88.3 \\
MATH500 & dflash2 & DFlash2-Unary-K16-B64 & 64 & 4.819 & 5.800 & 38.3 & 88.3 \\
MATH500 & original & Sequential-Baseline & -- & 0.000 & 1.000 & -- & 89.1 \\
MATH500 & original & DFlash-Original-DDTree-B16 & 16 & 6.981 & 7.981 & 43.0 & 87.5 \\
MATH500 & original & DFlash-Original-DDTree-B32 & 32 & 7.662 & 8.662 & 48.4 & 88.3 \\
MATH500 & original & DFlash-Original-DDTree-B64 & 64 & 8.177 & 9.177 & 43.8 & 89.1 \\
MATH500 & original & Raw-DFlash & -- & 6.341 & 7.341 & 46.1 & 87.5 \\
HumanEval & dflash2 & Sequential-Baseline & -- & 0.000 & 1.000 & -- & 94.5 \\
HumanEval & dflash2 & Raw-DFlash2 & -- & 3.227 & 4.214 & 43.9 & 93.9 \\
HumanEval & dflash2 & DFlash2-Original-DDTree-B16 & 16 & 3.819 & 4.804 & 42.7 & 93.9 \\
HumanEval & dflash2 & DFlash2-Original-DDTree-B32 & 32 & 4.126 & 5.109 & 43.9 & 94.5 \\
HumanEval & dflash2 & DFlash2-Original-DDTree-B64 & 64 & 4.357 & 5.340 & 47.0 & 93.9 \\
HumanEval & dflash2 & DFlash2-Original-DDTree-B7 & 7 & 3.347 & 4.333 & 40.2 & 94.5 \\
HumanEval & dflash2 & DFlash2-Pairwise-K16-B16 & 16 & 4.197 & 5.180 & 45.1 & 95.1 \\
HumanEval & dflash2 & DFlash2-Pairwise-K16-B32 & 32 & 4.577 & 5.560 & 49.4 & 93.9 \\
HumanEval & dflash2 & DFlash2-Pairwise-K16-B64 & 64 & 4.861 & 5.843 & 43.3 & 95.7 \\
HumanEval & dflash2 & DFlash2-Pairwise-K16-B7 & 7 & 3.542 & 4.528 & 37.8 & 95.1 \\
HumanEval & dflash2 & DFlash2-Unary-K16-B32 & 32 & 4.120 & 5.104 & 44.5 & 94.5 \\
HumanEval & dflash2 & DFlash2-Unary-K16-B64 & 64 & 4.360 & 5.343 & 48.2 & 93.9 \\
HumanEval & original & Sequential-Baseline & -- & 0.000 & 1.000 & -- & 94.5 \\
HumanEval & original & DFlash-Original-DDTree-B16 & 16 & 6.591 & 7.591 & 39.0 & 93.9 \\
HumanEval & original & DFlash-Original-DDTree-B32 & 32 & 7.288 & 8.288 & 44.5 & 94.5 \\
HumanEval & original & DFlash-Original-DDTree-B64 & 64 & 7.675 & 8.675 & 41.5 & 93.9 \\
HumanEval & original & Raw-DFlash & -- & 5.692 & 6.692 & 47.6 & 92.7 \\
MT-Bench-controlled & dflash2 & Sequential-Baseline & -- & 0.000 & 1.000 & -- & 63.7 \\
MT-Bench-controlled & dflash2 & Raw-DFlash2 & -- & 2.519 & 3.503 & 36.2 & 63.1 \\
MT-Bench-controlled & dflash2 & DFlash2-Original-DDTree-B16 & 16 & 3.116 & 4.097 & 36.2 & 64.4 \\
MT-Bench-controlled & dflash2 & DFlash2-Original-DDTree-B32 & 32 & 3.359 & 4.337 & 38.1 & 64.4 \\
MT-Bench-controlled & dflash2 & DFlash2-Original-DDTree-B64 & 64 & 3.585 & 4.565 & 40.0 & 63.7 \\
MT-Bench-controlled & dflash2 & DFlash2-Original-DDTree-B7 & 7 & 2.710 & 3.692 & 35.0 & 63.7 \\
MT-Bench-controlled & dflash2 & DFlash2-Pairwise-K16-B16 & 16 & 3.360 & 4.338 & 36.2 & 65.0 \\
MT-Bench-controlled & dflash2 & DFlash2-Pairwise-K16-B32 & 32 & 3.679 & 4.656 & 37.5 & 65.6 \\
MT-Bench-controlled & dflash2 & DFlash2-Pairwise-K16-B64 & 64 & 3.925 & 4.901 & 36.9 & 63.7 \\
MT-Bench-controlled & dflash2 & DFlash2-Pairwise-K16-B7 & 7 & 2.871 & 3.849 & 37.5 & 64.4 \\
MT-Bench-controlled & dflash2 & DFlash2-Unary-K16-B32 & 32 & 3.372 & 4.352 & 37.5 & 64.4 \\
MT-Bench-controlled & dflash2 & DFlash2-Unary-K16-B64 & 64 & 3.584 & 4.564 & 38.1 & 63.7 \\
MT-Bench-controlled & original & Sequential-Baseline & -- & 0.000 & 1.000 & -- & 63.7 \\
MT-Bench-controlled & original & DFlash-Original-DDTree-B16 & 16 & 4.061 & 5.061 & 33.8 & 63.1 \\
MT-Bench-controlled & original & DFlash-Original-DDTree-B32 & 32 & 4.480 & 5.480 & 39.4 & 63.1 \\
MT-Bench-controlled & original & DFlash-Original-DDTree-B64 & 64 & 4.837 & 5.837 & 36.9 & 65.0 \\
MT-Bench-controlled & original & Raw-DFlash & -- & 3.273 & 4.273 & 37.5 & 63.7 \\
MT-Bench-native & dflash2 & Sequential-Baseline & -- & 0.000 & 1.000 & -- & 63.7 \\
MT-Bench-native & dflash2 & Raw-DFlash2 & -- & 2.529 & 3.512 & 37.5 & 63.7 \\
MT-Bench-native & dflash2 & DFlash2-Original-DDTree-B16 & 16 & 3.109 & 4.091 & 36.2 & 63.7 \\
MT-Bench-native & dflash2 & DFlash2-Original-DDTree-B32 & 32 & 3.364 & 4.343 & 36.2 & 64.4 \\
MT-Bench-native & dflash2 & DFlash2-Original-DDTree-B64 & 64 & 3.591 & 4.571 & 45.0 & 63.1 \\
MT-Bench-native & dflash2 & DFlash2-Original-DDTree-B7 & 7 & 2.721 & 3.703 & 33.8 & 62.5 \\
MT-Bench-native & dflash2 & DFlash2-Pairwise-K16-B16 & 16 & 3.360 & 4.338 & 37.5 & 64.4 \\
MT-Bench-native & dflash2 & DFlash2-Pairwise-K16-B32 & 32 & 3.695 & 4.673 & 36.2 & 64.4 \\
MT-Bench-native & dflash2 & DFlash2-Pairwise-K16-B64 & 64 & 3.928 & 4.904 & 38.8 & 62.5 \\
MT-Bench-native & dflash2 & DFlash2-Pairwise-K16-B7 & 7 & 2.873 & 3.852 & 40.0 & 64.4 \\
MT-Bench-native & dflash2 & DFlash2-Unary-K16-B32 & 32 & 3.368 & 4.349 & 35.0 & 64.4 \\
MT-Bench-native & dflash2 & DFlash2-Unary-K16-B64 & 64 & 3.571 & 4.551 & 40.0 & 61.9 \\
MT-Bench-native & original & Sequential-Baseline & -- & 0.000 & 1.000 & -- & 63.7 \\
MT-Bench-native & original & DFlash-Original-DDTree-B16 & 16 & 4.066 & 5.066 & 36.2 & 61.9 \\
MT-Bench-native & original & DFlash-Original-DDTree-B32 & 32 & 4.484 & 5.484 & 42.5 & 63.7 \\
MT-Bench-native & original & DFlash-Original-DDTree-B64 & 64 & 4.867 & 5.867 & 36.2 & 64.4 \\
MT-Bench-native & original & Raw-DFlash & -- & 3.268 & 4.268 & 36.2 & 62.5 \\
\end{longtable}


\begin{longtable}{llllrrrrr}
\caption{All paired throughput comparisons. I/T/H denotes improved, tied, and hurt prompt clusters.}\label{tab:all-comparisons}\\
\toprule
Dataset & Type & Left & Right & B & $\Delta$ tok/s [95\% CI] & $\Delta$\% & Baseline-normalized [CI] & I/T/H \\
\midrule
\endfirsthead
\toprule
Dataset & Type & Left & Right & B & $\Delta$ tok/s [95\% CI] & $\Delta$\% & Baseline-normalized [CI] & I/T/H \\
\midrule
\endhead
\midrule
\multicolumn{9}{r}{Continued on next page}\\
\endfoot
\bottomrule
\endlastfoot
GSM8K & controlled & DFlash2-Pairwise-K16-B16 & DFlash2-Original-DDTree-B16 & 16 & 13.47 [12.11, 14.81] & 8.44 & -- & 121/0/7 \\
GSM8K & controlled & DFlash2-Pairwise-K16-B7 & DFlash2-Original-DDTree-B7 & 7 & 6.37 [4.88, 7.86] & 4.37 & -- & 102/0/26 \\
GSM8K & controlled & DFlash2-Pairwise-K16-B32 & DFlash2-Original-DDTree-B32 & 32 & 15.96 [14.41, 17.46] & 9.66 & -- & 124/0/4 \\
GSM8K & controlled & DFlash2-Pairwise-K16-B64 & DFlash2-Original-DDTree-B64 & 64 & 24.88 [23.09, 26.69] & 15.55 & -- & 128/0/0 \\
GSM8K & controlled & DFlash2-Pairwise-K16-B7 & Raw-DFlash2 & 7 & 2.61 [1.24, 3.98] & 1.89 & -- & 75/0/53 \\
GSM8K & controlled & DFlash2-Pairwise-K16-B32 & DFlash2-Unary-K16-B32 & 32 & 14.10 [12.59, 15.61] & 8.43 & -- & 122/0/6 \\
GSM8K & controlled & DFlash2-Pairwise-K16-B64 & DFlash2-Unary-K16-B64 & 64 & 15.26 [13.63, 16.93] & 9.04 & -- & 123/0/5 \\
GSM8K & cross\_drafter & DFlash2-Pairwise-K16-B16 & Raw-DFlash & 16 & -56.18 [-61.32, -51.13] & -22.77 & -0.832 [-0.927, -0.737] & 3/0/125 \\
GSM8K & cross\_drafter & DFlash2-Pairwise-K16-B16 & DFlash-Original-DDTree-B16 & 16 & -84.43 [-89.43, -79.56] & -31.54 & -1.365 [-1.457, -1.277] & 0/0/128 \\
GSM8K & cross\_drafter & DFlash2-Pairwise-K16-B32 & Raw-DFlash & 32 & -47.73 [-52.87, -42.53] & -19.04 & -0.662 [-0.759, -0.570] & 7/0/121 \\
GSM8K & cross\_drafter & DFlash2-Pairwise-K16-B32 & DFlash-Original-DDTree-B32 & 32 & -93.71 [-99.12, -88.40] & -32.75 & -1.528 [-1.629, -1.427] & 0/0/128 \\
GSM8K & cross\_drafter & DFlash2-Pairwise-K16-B64 & Raw-DFlash & 64 & -43.95 [-49.28, -38.63] & -17.32 & -0.586 [-0.688, -0.488] & 6/0/122 \\
GSM8K & cross\_drafter & DFlash2-Pairwise-K16-B64 & DFlash-Original-DDTree-B64 & 64 & -105.43 [-111.08, -99.87] & -35.06 & -1.744 [-1.852, -1.641] & 0/0/128 \\
GSM8K & within\_family\_vs\_sequential & Raw-DFlash & Sequential-Baseline & -- & 179.90 [173.14, 186.78] & 338.93 & -- & 128/0/0 \\
GSM8K & within\_family\_vs\_sequential & DFlash-Original-DDTree-B16 & Sequential-Baseline & 16 & 208.15 [201.68, 214.81] & 392.19 & -- & 128/0/0 \\
GSM8K & within\_family\_vs\_sequential & DFlash-Original-DDTree-B32 & Sequential-Baseline & 32 & 225.87 [218.81, 232.86] & 425.54 & -- & 128/0/0 \\
GSM8K & within\_family\_vs\_sequential & DFlash-Original-DDTree-B64 & Sequential-Baseline & 64 & 241.38 [234.15, 248.62] & 454.75 & -- & 128/0/0 \\
GSM8K & controlled & DFlash-Original-DDTree-B16 & Raw-DFlash & 16 & 28.25 [25.61, 30.87] & 12.79 & -- & 127/0/1 \\
GSM8K & controlled & DFlash-Original-DDTree-B32 & Raw-DFlash & 32 & 45.98 [43.15, 48.74] & 20.40 & -- & 128/0/0 \\
GSM8K & controlled & DFlash-Original-DDTree-B64 & Raw-DFlash & 64 & 61.48 [58.30, 64.68] & 27.31 & -- & 127/0/1 \\
GSM8K & within\_family\_vs\_sequential & Raw-DFlash2 & Sequential-Baseline & -- & 106.04 [103.09, 108.98] & 213.42 & -- & 128/0/0 \\
GSM8K & within\_family\_vs\_sequential & DFlash2-Original-DDTree-B7 & Sequential-Baseline & 7 & 102.28 [99.49, 105.02] & 205.76 & -- & 128/0/0 \\
GSM8K & within\_family\_vs\_sequential & DFlash2-Original-DDTree-B16 & Sequential-Baseline & 16 & 113.62 [111.02, 116.25] & 228.59 & -- & 128/0/0 \\
GSM8K & within\_family\_vs\_sequential & DFlash2-Original-DDTree-B32 & Sequential-Baseline & 32 & 119.57 [117.01, 122.14] & 240.56 & -- & 128/0/0 \\
GSM8K & within\_family\_vs\_sequential & DFlash2-Original-DDTree-B64 & Sequential-Baseline & 64 & 114.44 [111.74, 117.03] & 230.20 & -- & 128/0/0 \\
GSM8K & within\_family\_vs\_sequential & DFlash2-Pairwise-K16-B7 & Sequential-Baseline & 7 & 108.65 [105.74, 111.56] & 218.60 & -- & 128/0/0 \\
GSM8K & within\_family\_vs\_sequential & DFlash2-Pairwise-K16-B16 & Sequential-Baseline & 16 & 127.09 [124.53, 129.69] & 255.71 & -- & 128/0/0 \\
GSM8K & within\_family\_vs\_sequential & DFlash2-Pairwise-K16-B32 & Sequential-Baseline & 32 & 135.54 [132.93, 138.15] & 272.73 & -- & 128/0/0 \\
GSM8K & within\_family\_vs\_sequential & DFlash2-Pairwise-K16-B64 & Sequential-Baseline & 64 & 139.32 [136.75, 141.74] & 280.31 & -- & 128/0/0 \\
GSM8K & within\_family\_vs\_sequential & DFlash2-Unary-K16-B32 & Sequential-Baseline & 32 & 121.44 [118.94, 124.03] & 244.34 & -- & 128/0/0 \\
GSM8K & within\_family\_vs\_sequential & DFlash2-Unary-K16-B64 & Sequential-Baseline & 64 & 124.06 [121.42, 126.65] & 249.56 & -- & 128/0/0 \\
HumanEval & controlled & DFlash2-Pairwise-K16-B16 & DFlash2-Original-DDTree-B16 & 16 & 10.66 [9.08, 12.28] & 7.61 & -- & 144/0/20 \\
HumanEval & controlled & DFlash2-Pairwise-K16-B7 & DFlash2-Original-DDTree-B7 & 7 & 5.62 [4.54, 6.69] & 4.39 & -- & 129/0/35 \\
HumanEval & controlled & DFlash2-Pairwise-K16-B32 & DFlash2-Original-DDTree-B32 & 32 & 14.29 [13.01, 15.55] & 9.75 & -- & 159/0/5 \\
HumanEval & controlled & DFlash2-Pairwise-K16-B64 & DFlash2-Original-DDTree-B64 & 64 & 23.31 [21.84, 24.81] & 16.33 & -- & 163/0/1 \\
HumanEval & controlled & DFlash2-Pairwise-K16-B7 & Raw-DFlash2 & 7 & 4.01 [2.92, 5.10] & 3.34 & -- & 121/0/43 \\
HumanEval & controlled & DFlash2-Pairwise-K16-B32 & DFlash2-Unary-K16-B32 & 32 & 12.94 [11.61, 14.26] & 8.79 & -- & 154/0/10 \\
HumanEval & controlled & DFlash2-Pairwise-K16-B64 & DFlash2-Unary-K16-B64 & 64 & 13.89 [12.50, 15.29] & 9.16 & -- & 156/0/8 \\
HumanEval & cross\_drafter & DFlash2-Pairwise-K16-B16 & Raw-DFlash & 16 & -84.16 [-91.13, -77.14] & -32.74 & -1.448 [-1.578, -1.320] & 2/0/162 \\
HumanEval & cross\_drafter & DFlash2-Pairwise-K16-B16 & DFlash-Original-DDTree-B16 & 16 & -110.43 [-116.72, -104.17] & -40.09 & -1.955 [-2.072, -1.835] & 0/0/164 \\
HumanEval & cross\_drafter & DFlash2-Pairwise-K16-B32 & Raw-DFlash & 32 & -75.49 [-82.51, -68.55] & -28.83 & -1.272 [-1.404, -1.143] & 2/0/162 \\
HumanEval & cross\_drafter & DFlash2-Pairwise-K16-B32 & DFlash-Original-DDTree-B32 & 32 & -123.26 [-130.13, -116.46] & -41.44 & -2.192 [-2.323, -2.065] & 0/0/164 \\
HumanEval & cross\_drafter & DFlash2-Pairwise-K16-B64 & Raw-DFlash & 64 & -69.53 [-76.81, -62.52] & -26.10 & -1.151 [-1.287, -1.018] & 5/0/159 \\
HumanEval & cross\_drafter & DFlash2-Pairwise-K16-B64 & DFlash-Original-DDTree-B64 & 64 & -131.61 [-138.24, -125.09] & -42.39 & -2.346 [-2.471, -2.223] & 0/0/164 \\
HumanEval & within\_family\_vs\_sequential & Raw-DFlash & Sequential-Baseline & -- & 188.69 [180.24, 197.50] & 363.38 & -- & 164/0/0 \\
HumanEval & within\_family\_vs\_sequential & DFlash-Original-DDTree-B16 & Sequential-Baseline & 16 & 214.96 [206.62, 223.17] & 414.01 & -- & 164/0/0 \\
HumanEval & within\_family\_vs\_sequential & DFlash-Original-DDTree-B32 & Sequential-Baseline & 32 & 236.46 [227.87, 245.40] & 455.37 & -- & 164/0/0 \\
HumanEval & within\_family\_vs\_sequential & DFlash-Original-DDTree-B64 & Sequential-Baseline & 64 & 250.77 [242.45, 259.30] & 482.93 & -- & 164/0/0 \\
HumanEval & controlled & DFlash-Original-DDTree-B16 & Raw-DFlash & 16 & 26.27 [23.69, 28.89] & 12.15 & -- & 154/0/10 \\
HumanEval & controlled & DFlash-Original-DDTree-B32 & Raw-DFlash & 32 & 47.78 [44.59, 50.87] & 21.37 & -- & 161/0/3 \\
HumanEval & controlled & DFlash-Original-DDTree-B64 & Raw-DFlash & 64 & 62.08 [57.77, 66.17] & 28.05 & -- & 160/0/4 \\
HumanEval & within\_family\_vs\_sequential & Raw-DFlash2 & Sequential-Baseline & -- & 86.32 [83.23, 89.48] & 175.74 & -- & 164/0/0 \\
HumanEval & within\_family\_vs\_sequential & DFlash2-Original-DDTree-B7 & Sequential-Baseline & 7 & 84.72 [81.80, 87.73] & 172.51 & -- & 164/0/0 \\
HumanEval & within\_family\_vs\_sequential & DFlash2-Original-DDTree-B16 & Sequential-Baseline & 16 & 96.66 [93.86, 99.56] & 196.84 & -- & 164/0/0 \\
HumanEval & within\_family\_vs\_sequential & DFlash2-Original-DDTree-B32 & Sequential-Baseline & 32 & 101.70 [98.92, 104.41] & 207.10 & -- & 164/0/0 \\
HumanEval & within\_family\_vs\_sequential & DFlash2-Original-DDTree-B64 & Sequential-Baseline & 64 & 98.65 [95.92, 101.44] & 200.78 & -- & 164/0/0 \\
HumanEval & within\_family\_vs\_sequential & DFlash2-Pairwise-K16-B7 & Sequential-Baseline & 7 & 90.33 [87.44, 93.29] & 183.93 & -- & 164/0/0 \\
HumanEval & within\_family\_vs\_sequential & DFlash2-Pairwise-K16-B16 & Sequential-Baseline & 16 & 107.32 [104.34, 110.36] & 218.55 & -- & 164/0/0 \\
HumanEval & within\_family\_vs\_sequential & DFlash2-Pairwise-K16-B32 & Sequential-Baseline & 32 & 116.00 [113.17, 118.82] & 236.21 & -- & 164/0/0 \\
HumanEval & within\_family\_vs\_sequential & DFlash2-Pairwise-K16-B64 & Sequential-Baseline & 64 & 121.96 [119.25, 124.67] & 248.29 & -- & 164/0/0 \\
HumanEval & within\_family\_vs\_sequential & DFlash2-Unary-K16-B32 & Sequential-Baseline & 32 & 103.06 [100.26, 105.89] & 209.86 & -- & 164/0/0 \\
HumanEval & within\_family\_vs\_sequential & DFlash2-Unary-K16-B64 & Sequential-Baseline & 64 & 108.07 [105.34, 110.83] & 219.99 & -- & 164/0/0 \\
MATH500 & controlled & DFlash2-Pairwise-K16-B16 & DFlash2-Original-DDTree-B16 & 16 & 15.63 [13.88, 17.42] & 9.60 & -- & 120/0/8 \\
MATH500 & controlled & DFlash2-Pairwise-K16-B7 & DFlash2-Original-DDTree-B7 & 7 & 8.88 [7.38, 10.44] & 5.94 & -- & 110/0/18 \\
MATH500 & controlled & DFlash2-Pairwise-K16-B32 & DFlash2-Original-DDTree-B32 & 32 & 19.30 [17.54, 21.09] & 11.40 & -- & 125/0/3 \\
MATH500 & controlled & DFlash2-Pairwise-K16-B64 & DFlash2-Original-DDTree-B64 & 64 & 28.28 [26.48, 30.06] & 17.33 & -- & 128/0/0 \\
MATH500 & controlled & DFlash2-Pairwise-K16-B7 & Raw-DFlash2 & 7 & 2.75 [1.23, 4.30] & 2.18 & -- & 81/0/47 \\
MATH500 & controlled & DFlash2-Pairwise-K16-B32 & DFlash2-Unary-K16-B32 & 32 & 17.73 [16.03, 19.47] & 10.44 & -- & 126/0/2 \\
MATH500 & controlled & DFlash2-Pairwise-K16-B64 & DFlash2-Unary-K16-B64 & 64 & 18.51 [16.88, 20.14] & 10.67 & -- & 124/0/4 \\
MATH500 & cross\_drafter & DFlash2-Pairwise-K16-B16 & Raw-DFlash & 16 & -84.06 [-93.98, -74.38] & -28.06 & -1.410 [-1.595, -1.233] & 2/0/126 \\
MATH500 & cross\_drafter & DFlash2-Pairwise-K16-B16 & DFlash-Original-DDTree-B16 & 16 & -101.17 [-109.83, -92.79] & -33.50 & -1.731 [-1.885, -1.581] & 1/0/127 \\
MATH500 & cross\_drafter & DFlash2-Pairwise-K16-B32 & Raw-DFlash & 32 & -76.35 [-86.53, -66.40] & -24.98 & -1.258 [-1.445, -1.079] & 6/0/122 \\
MATH500 & cross\_drafter & DFlash2-Pairwise-K16-B32 & DFlash-Original-DDTree-B32 & 32 & -115.55 [-123.90, -107.22] & -36.00 & -1.994 [-2.152, -1.842] & 0/0/128 \\
MATH500 & cross\_drafter & DFlash2-Pairwise-K16-B64 & Raw-DFlash & 64 & -72.09 [-82.23, -62.00] & -23.11 & -1.175 [-1.364, -0.991] & 12/0/116 \\
MATH500 & cross\_drafter & DFlash2-Pairwise-K16-B64 & DFlash-Original-DDTree-B64 & 64 & -130.37 [-140.11, -120.83] & -38.23 & -2.268 [-2.442, -2.099] & 0/0/128 \\
MATH500 & within\_family\_vs\_sequential & Raw-DFlash & Sequential-Baseline & -- & 214.90 [202.53, 227.62] & 402.18 & -- & 128/0/0 \\
MATH500 & within\_family\_vs\_sequential & DFlash-Original-DDTree-B16 & Sequential-Baseline & 16 & 232.01 [221.09, 243.34] & 434.35 & -- & 128/0/0 \\
MATH500 & within\_family\_vs\_sequential & DFlash-Original-DDTree-B32 & Sequential-Baseline & 32 & 254.11 [243.30, 265.35] & 475.75 & -- & 128/0/0 \\
MATH500 & within\_family\_vs\_sequential & DFlash-Original-DDTree-B64 & Sequential-Baseline & 64 & 273.19 [261.45, 285.01] & 511.47 & -- & 128/0/0 \\
MATH500 & controlled & DFlash-Original-DDTree-B16 & Raw-DFlash & 16 & 17.11 [12.94, 21.13] & 7.88 & -- & 106/0/22 \\
MATH500 & controlled & DFlash-Original-DDTree-B32 & Raw-DFlash & 32 & 39.20 [34.23, 44.20] & 16.86 & -- & 114/0/14 \\
MATH500 & controlled & DFlash-Original-DDTree-B64 & Raw-DFlash & 64 & 58.28 [53.91, 62.56] & 23.97 & -- & 125/0/3 \\
MATH500 & within\_family\_vs\_sequential & Raw-DFlash2 & Sequential-Baseline & -- & 111.26 [106.84, 115.85] & 218.11 & -- & 128/0/0 \\
MATH500 & within\_family\_vs\_sequential & DFlash2-Original-DDTree-B7 & Sequential-Baseline & 7 & 105.14 [101.31, 109.03] & 206.12 & -- & 128/0/0 \\
MATH500 & within\_family\_vs\_sequential & DFlash2-Original-DDTree-B16 & Sequential-Baseline & 16 & 117.60 [113.91, 121.24] & 230.57 & -- & 128/0/0 \\
MATH500 & within\_family\_vs\_sequential & DFlash2-Original-DDTree-B32 & Sequential-Baseline & 32 & 121.65 [118.15, 125.06] & 238.56 & -- & 128/0/0 \\
MATH500 & within\_family\_vs\_sequential & DFlash2-Original-DDTree-B64 & Sequential-Baseline & 64 & 116.92 [113.49, 120.35] & 229.27 & -- & 128/0/0 \\
MATH500 & within\_family\_vs\_sequential & DFlash2-Pairwise-K16-B7 & Sequential-Baseline & 7 & 114.01 [110.12, 117.94] & 223.54 & -- & 128/0/0 \\
MATH500 & within\_family\_vs\_sequential & DFlash2-Pairwise-K16-B16 & Sequential-Baseline & 16 & 133.23 [129.52, 136.93] & 261.21 & -- & 128/0/0 \\
MATH500 & within\_family\_vs\_sequential & DFlash2-Pairwise-K16-B32 & Sequential-Baseline & 32 & 140.95 [137.23, 144.55] & 276.37 & -- & 128/0/0 \\
MATH500 & within\_family\_vs\_sequential & DFlash2-Pairwise-K16-B64 & Sequential-Baseline & 64 & 145.20 [141.78, 148.50] & 284.65 & -- & 128/0/0 \\
MATH500 & within\_family\_vs\_sequential & DFlash2-Unary-K16-B32 & Sequential-Baseline & 32 & 123.21 [119.64, 126.78] & 241.62 & -- & 128/0/0 \\
MATH500 & within\_family\_vs\_sequential & DFlash2-Unary-K16-B64 & Sequential-Baseline & 64 & 126.69 [123.36, 129.95] & 248.43 & -- & 128/0/0 \\
MT-Bench-controlled & controlled & DFlash2-Pairwise-K16-B16 & DFlash2-Original-DDTree-B16 & 16 & 7.07 [5.53, 8.62] & 5.41 & -- & 67/0/13 \\
MT-Bench-controlled & controlled & DFlash2-Pairwise-K16-B7 & DFlash2-Original-DDTree-B7 & 7 & 4.65 [3.40, 5.94] & 4.00 & -- & 64/0/16 \\
MT-Bench-controlled & controlled & DFlash2-Pairwise-K16-B32 & DFlash2-Original-DDTree-B32 & 32 & 10.56 [8.81, 12.34] & 8.11 & -- & 75/0/5 \\
MT-Bench-controlled & controlled & DFlash2-Pairwise-K16-B64 & DFlash2-Original-DDTree-B64 & 64 & 18.13 [15.38, 20.56] & 14.35 & -- & 78/0/2 \\
MT-Bench-controlled & controlled & DFlash2-Pairwise-K16-B7 & Raw-DFlash2 & 7 & 5.88 [4.53, 7.24] & 6.41 & -- & 66/0/14 \\
MT-Bench-controlled & controlled & DFlash2-Pairwise-K16-B32 & DFlash2-Unary-K16-B32 & 32 & 8.70 [6.75, 10.52] & 6.61 & -- & 72/0/8 \\
MT-Bench-controlled & controlled & DFlash2-Pairwise-K16-B64 & DFlash2-Unary-K16-B64 & 64 & 8.85 [6.39, 11.05] & 6.48 & -- & 70/0/10 \\
MT-Bench-controlled & cross\_drafter & DFlash2-Pairwise-K16-B16 & Raw-DFlash & 16 & -21.65 [-31.80, -12.31] & -5.21 & -0.283 [-0.462, -0.120] & 39/0/41 \\
MT-Bench-controlled & cross\_drafter & DFlash2-Pairwise-K16-B16 & DFlash-Original-DDTree-B16 & 16 & -46.93 [-56.46, -38.02] & -21.76 & -0.753 [-0.920, -0.604] & 0/0/80 \\
MT-Bench-controlled & cross\_drafter & DFlash2-Pairwise-K16-B32 & Raw-DFlash & 32 & -13.67 [-24.10, -4.11] & 1.16 & -0.129 [-0.311, 0.036] & 48/0/32 \\
MT-Bench-controlled & cross\_drafter & DFlash2-Pairwise-K16-B32 & DFlash-Original-DDTree-B32 & 32 & -52.56 [-63.14, -42.91] & -22.46 & -0.851 [-1.041, -0.679] & 0/0/80 \\
MT-Bench-controlled & cross\_drafter & DFlash2-Pairwise-K16-B64 & Raw-DFlash & 64 & -9.46 [-19.67, 0.01] & 4.47 & -0.047 [-0.227, 0.123] & 47/0/33 \\
MT-Bench-controlled & cross\_drafter & DFlash2-Pairwise-K16-B64 & DFlash-Original-DDTree-B64 & 64 & -61.24 [-72.52, -50.45] & -25.05 & -1.009 [-1.214, -0.819] & 0/0/80 \\
MT-Bench-controlled & within\_family\_vs\_sequential & Raw-DFlash & Sequential-Baseline & -- & 104.48 [88.04, 122.54] & 193.96 & -- & 80/0/0 \\
MT-Bench-controlled & within\_family\_vs\_sequential & DFlash-Original-DDTree-B16 & Sequential-Baseline & 16 & 129.77 [113.64, 146.94] & 240.95 & -- & 80/0/0 \\
MT-Bench-controlled & within\_family\_vs\_sequential & DFlash-Original-DDTree-B32 & Sequential-Baseline & 32 & 143.37 [126.05, 161.74] & 266.19 & -- & 80/0/0 \\
MT-Bench-controlled & within\_family\_vs\_sequential & DFlash-Original-DDTree-B64 & Sequential-Baseline & 64 & 156.25 [138.16, 175.36] & 290.14 & -- & 80/0/0 \\
MT-Bench-controlled & controlled & DFlash-Original-DDTree-B16 & Raw-DFlash & 16 & 25.29 [22.96, 27.54] & 20.35 & -- & 76/0/4 \\
MT-Bench-controlled & controlled & DFlash-Original-DDTree-B32 & Raw-DFlash & 32 & 38.89 [36.05, 41.68] & 29.41 & -- & 79/0/1 \\
MT-Bench-controlled & controlled & DFlash-Original-DDTree-B64 & Raw-DFlash & 64 & 51.77 [48.23, 55.57] & 38.28 & -- & 80/0/0 \\
MT-Bench-controlled & within\_family\_vs\_sequential & Raw-DFlash2 & Sequential-Baseline & -- & 66.28 [57.97, 74.96] & 129.07 & -- & 80/0/0 \\
MT-Bench-controlled & within\_family\_vs\_sequential & DFlash2-Original-DDTree-B7 & Sequential-Baseline & 7 & 67.50 [60.11, 75.28] & 131.38 & -- & 80/0/0 \\
MT-Bench-controlled & within\_family\_vs\_sequential & DFlash2-Original-DDTree-B16 & Sequential-Baseline & 16 & 78.05 [70.34, 85.96] & 151.91 & -- & 80/0/0 \\
MT-Bench-controlled & within\_family\_vs\_sequential & DFlash2-Original-DDTree-B32 & Sequential-Baseline & 32 & 82.52 [74.97, 90.35] & 160.56 & -- & 80/0/0 \\
MT-Bench-controlled & within\_family\_vs\_sequential & DFlash2-Original-DDTree-B64 & Sequential-Baseline & 64 & 79.17 [71.81, 86.89] & 154.05 & -- & 80/0/0 \\
MT-Bench-controlled & within\_family\_vs\_sequential & DFlash2-Pairwise-K16-B7 & Sequential-Baseline & 7 & 72.16 [64.43, 79.91] & 140.46 & -- & 80/0/0 \\
MT-Bench-controlled & within\_family\_vs\_sequential & DFlash2-Pairwise-K16-B16 & Sequential-Baseline & 16 & 85.11 [76.88, 93.56] & 165.63 & -- & 80/0/0 \\
MT-Bench-controlled & within\_family\_vs\_sequential & DFlash2-Pairwise-K16-B32 & Sequential-Baseline & 32 & 93.09 [84.95, 101.49] & 181.08 & -- & 80/0/0 \\
MT-Bench-controlled & within\_family\_vs\_sequential & DFlash2-Pairwise-K16-B64 & Sequential-Baseline & 64 & 97.30 [89.21, 105.89] & 189.26 & -- & 80/0/0 \\
MT-Bench-controlled & within\_family\_vs\_sequential & DFlash2-Unary-K16-B32 & Sequential-Baseline & 32 & 84.39 [76.63, 92.36] & 164.21 & -- & 80/0/0 \\
MT-Bench-controlled & within\_family\_vs\_sequential & DFlash2-Unary-K16-B64 & Sequential-Baseline & 64 & 88.45 [80.80, 96.30] & 172.07 & -- & 80/0/0 \\
MT-Bench-native & controlled & DFlash2-Pairwise-K16-B16 & DFlash2-Original-DDTree-B16 & 16 & 7.20 [5.43, 9.02] & 5.32 & -- & 68/0/12 \\
MT-Bench-native & controlled & DFlash2-Pairwise-K16-B7 & DFlash2-Original-DDTree-B7 & 7 & 4.51 [3.09, 6.00] & 3.82 & -- & 62/0/18 \\
MT-Bench-native & controlled & DFlash2-Pairwise-K16-B32 & DFlash2-Original-DDTree-B32 & 32 & 10.68 [8.87, 12.49] & 8.10 & -- & 73/0/7 \\
MT-Bench-native & controlled & DFlash2-Pairwise-K16-B64 & DFlash2-Original-DDTree-B64 & 64 & 18.70 [16.40, 20.97] & 14.49 & -- & 76/0/4 \\
MT-Bench-native & controlled & DFlash2-Pairwise-K16-B7 & Raw-DFlash2 & 7 & 5.68 [4.24, 7.08] & 6.23 & -- & 64/0/16 \\
MT-Bench-native & controlled & DFlash2-Pairwise-K16-B32 & DFlash2-Unary-K16-B32 & 32 & 9.31 [7.15, 11.39] & 7.08 & -- & 71/0/9 \\
MT-Bench-native & controlled & DFlash2-Pairwise-K16-B64 & DFlash2-Unary-K16-B64 & 64 & 10.38 [8.34, 12.44] & 7.45 & -- & 72/0/8 \\
MT-Bench-native & cross\_drafter & DFlash2-Pairwise-K16-B16 & Raw-DFlash & 16 & -15.32 [-24.46, -6.60] & -1.83 & -0.280 [-0.454, -0.121] & 44/0/36 \\
MT-Bench-native & cross\_drafter & DFlash2-Pairwise-K16-B16 & DFlash-Original-DDTree-B16 & 16 & -40.29 [-49.33, -31.82] & -18.64 & -0.759 [-0.930, -0.601] & 0/0/80 \\
MT-Bench-native & cross\_drafter & DFlash2-Pairwise-K16-B32 & Raw-DFlash & 32 & -6.95 [-16.90, 1.90] & 5.37 & -0.119 [-0.296, 0.045] & 50/0/30 \\
MT-Bench-native & cross\_drafter & DFlash2-Pairwise-K16-B32 & DFlash-Original-DDTree-B32 & 32 & -44.99 [-55.28, -35.61] & -19.17 & -0.849 [-1.034, -0.674] & 1/0/79 \\
MT-Bench-native & cross\_drafter & DFlash2-Pairwise-K16-B64 & Raw-DFlash & 64 & -2.04 [-11.69, 6.77] & 9.12 & -0.024 [-0.198, 0.138] & 53/0/27 \\
MT-Bench-native & cross\_drafter & DFlash2-Pairwise-K16-B64 & DFlash-Original-DDTree-B64 & 64 & -53.30 [-64.21, -43.29] & -21.78 & -1.008 [-1.216, -0.815] & 1/0/79 \\
MT-Bench-native & within\_family\_vs\_sequential & Raw-DFlash & Sequential-Baseline & -- & 100.90 [84.90, 117.61] & 193.69 & -- & 80/0/0 \\
MT-Bench-native & within\_family\_vs\_sequential & DFlash-Original-DDTree-B16 & Sequential-Baseline & 16 & 125.88 [110.30, 142.48] & 241.65 & -- & 80/0/0 \\
MT-Bench-native & within\_family\_vs\_sequential & DFlash-Original-DDTree-B32 & Sequential-Baseline & 32 & 138.94 [122.31, 156.45] & 266.74 & -- & 80/0/0 \\
MT-Bench-native & within\_family\_vs\_sequential & DFlash-Original-DDTree-B64 & Sequential-Baseline & 64 & 152.16 [134.48, 170.89] & 292.13 & -- & 80/0/0 \\
MT-Bench-native & controlled & DFlash-Original-DDTree-B16 & Raw-DFlash & 16 & 24.98 [22.55, 27.39] & 20.02 & -- & 77/0/3 \\
MT-Bench-native & controlled & DFlash-Original-DDTree-B32 & Raw-DFlash & 32 & 38.04 [34.90, 41.11] & 29.39 & -- & 78/0/2 \\
MT-Bench-native & controlled & DFlash-Original-DDTree-B64 & Raw-DFlash & 64 & 51.26 [47.19, 55.38] & 38.65 & -- & 80/0/0 \\
MT-Bench-native & within\_family\_vs\_sequential & Raw-DFlash2 & Sequential-Baseline & -- & 67.27 [58.59, 76.17] & 130.08 & -- & 80/0/0 \\
MT-Bench-native & within\_family\_vs\_sequential & DFlash2-Original-DDTree-B7 & Sequential-Baseline & 7 & 68.44 [60.78, 76.42] & 132.28 & -- & 80/0/0 \\
MT-Bench-native & within\_family\_vs\_sequential & DFlash2-Original-DDTree-B16 & Sequential-Baseline & 16 & 78.55 [70.88, 86.30] & 151.80 & -- & 80/0/0 \\
MT-Bench-native & within\_family\_vs\_sequential & DFlash2-Original-DDTree-B32 & Sequential-Baseline & 32 & 83.44 [75.81, 91.23] & 161.22 & -- & 80/0/0 \\
MT-Bench-native & within\_family\_vs\_sequential & DFlash2-Original-DDTree-B64 & Sequential-Baseline & 64 & 80.33 [73.03, 87.82] & 155.18 & -- & 80/0/0 \\
MT-Bench-native & within\_family\_vs\_sequential & DFlash2-Pairwise-K16-B7 & Sequential-Baseline & 7 & 72.95 [65.06, 81.30] & 141.00 & -- & 80/0/0 \\
MT-Bench-native & within\_family\_vs\_sequential & DFlash2-Pairwise-K16-B16 & Sequential-Baseline & 16 & 85.76 [77.42, 94.37] & 165.73 & -- & 80/0/0 \\
MT-Bench-native & within\_family\_vs\_sequential & DFlash2-Pairwise-K16-B32 & Sequential-Baseline & 32 & 94.12 [86.04, 102.17] & 181.84 & -- & 80/0/0 \\
MT-Bench-native & within\_family\_vs\_sequential & DFlash2-Pairwise-K16-B64 & Sequential-Baseline & 64 & 99.03 [90.85, 107.54] & 191.32 & -- & 80/0/0 \\
MT-Bench-native & within\_family\_vs\_sequential & DFlash2-Unary-K16-B32 & Sequential-Baseline & 32 & 84.81 [76.99, 92.71] & 163.86 & -- & 80/0/0 \\
MT-Bench-native & within\_family\_vs\_sequential & DFlash2-Unary-K16-B64 & Sequential-Baseline & 64 & 88.66 [81.17, 96.64] & 171.26 & -- & 80/0/0 \\
\end{longtable}


\begin{longtable}{lllrrrrrrrrr}
\caption{All timing decompositions. Stage values are milliseconds per call.}\label{tab:all-timing}\\
\toprule
Dataset & Left & Right & $\Delta$ match & $\Delta$ commit & Calls saved & Draft L & Draft R & Tree & Verify L & Verify R & $\Delta$ tok/s \\
\midrule
\endfirsthead
\toprule
Dataset & Left & Right & $\Delta$ match & $\Delta$ commit & Calls saved & Draft L & Draft R & Tree & Verify L & Verify R & $\Delta$ tok/s \\
\midrule
\endhead
\midrule
\multicolumn{12}{r}{Continued on next page}\\
\endfoot
\bottomrule
\endlastfoot
GSM8K & DFlash2-Pairwise-K16-B16 & DFlash2-Original-DDTree-B16 & 0.476 & 0.473 & 3.61 & 8.093 & 8.108 & 0.125 & 21.879 & 21.865 & 13.47 \\
GSM8K & DFlash2-Pairwise-K16-B7 & DFlash2-Original-DDTree-B7 & 0.224 & 0.221 & 2.20 & 8.085 & 8.078 & 0.096 & 21.652 & 21.647 & 6.37 \\
GSM8K & DFlash2-Pairwise-K16-B32 & DFlash2-Original-DDTree-B32 & 0.519 & 0.517 & 3.62 & 8.093 & 8.106 & -0.146 & 21.960 & 21.953 & 15.96 \\
GSM8K & DFlash2-Pairwise-K16-B64 & DFlash2-Original-DDTree-B64 & 0.552 & 0.550 & 3.61 & 8.085 & 8.121 & -1.812 & 21.833 & 21.866 & 24.88 \\
GSM8K & DFlash2-Pairwise-K16-B7 & Raw-DFlash2 & 0.290 & 0.288 & 3.09 & 8.085 & 7.319 & 0.785 & 21.652 & 22.514 & 2.61 \\
GSM8K & DFlash2-Pairwise-K16-B32 & DFlash2-Unary-K16-B32 & 0.512 & 0.510 & 3.56 & 8.093 & 8.094 & 0.138 & 21.960 & 21.914 & 14.10 \\
GSM8K & DFlash2-Pairwise-K16-B64 & DFlash2-Unary-K16-B64 & 0.559 & 0.557 & 3.69 & 8.085 & 8.099 & 0.158 & 21.833 & 21.808 & 15.26 \\
GSM8K & DFlash2-Pairwise-K16-B16 & Raw-DFlash & -0.520 & -0.542 & -1.97 & 8.093 & 4.514 & 0.958 & 21.879 & 21.211 & -56.18 \\
GSM8K & DFlash2-Pairwise-K16-B16 & DFlash-Original-DDTree-B16 & -1.438 & -1.460 & -7.20 & 8.093 & 4.515 & 0.471 & 21.879 & 20.306 & -84.43 \\
GSM8K & DFlash2-Pairwise-K16-B32 & Raw-DFlash & -0.164 & -0.187 & 0.48 & 8.093 & 4.514 & 1.255 & 21.960 & 21.211 & -47.73 \\
GSM8K & DFlash2-Pairwise-K16-B32 & DFlash-Original-DDTree-B32 & -1.642 & -1.666 & -7.09 & 8.093 & 4.517 & 0.729 & 21.960 & 20.402 & -93.71 \\
GSM8K & DFlash2-Pairwise-K16-B64 & Raw-DFlash & 0.053 & 0.030 & 1.86 & 8.085 & 4.514 & 1.842 & 21.833 & 21.211 & -43.95 \\
GSM8K & DFlash2-Pairwise-K16-B64 & DFlash-Original-DDTree-B64 & -1.833 & -1.857 & -7.38 & 8.085 & 4.504 & 1.249 & 21.833 & 20.202 & -105.43 \\
HumanEval & DFlash2-Pairwise-K16-B16 & DFlash2-Original-DDTree-B16 & 0.378 & 0.377 & 3.81 & 8.235 & 8.251 & 0.136 & 22.094 & 22.090 & 10.66 \\
HumanEval & DFlash2-Pairwise-K16-B7 & DFlash2-Original-DDTree-B7 & 0.195 & 0.195 & 2.57 & 8.231 & 8.237 & 0.110 & 21.911 & 21.936 & 5.62 \\
HumanEval & DFlash2-Pairwise-K16-B32 & DFlash2-Original-DDTree-B32 & 0.451 & 0.451 & 4.18 & 8.246 & 8.277 & -0.181 & 22.271 & 22.321 & 14.29 \\
HumanEval & DFlash2-Pairwise-K16-B64 & DFlash2-Original-DDTree-B64 & 0.504 & 0.503 & 4.13 & 8.231 & 8.264 & -2.065 & 22.122 & 22.152 & 23.31 \\
HumanEval & DFlash2-Pairwise-K16-B7 & Raw-DFlash2 & 0.315 & 0.314 & 4.48 & 8.231 & 7.461 & 0.815 & 21.911 & 22.834 & 4.01 \\
HumanEval & DFlash2-Pairwise-K16-B32 & DFlash2-Unary-K16-B32 & 0.457 & 0.456 & 4.25 & 8.246 & 8.275 & 0.140 & 22.271 & 22.302 & 12.94 \\
HumanEval & DFlash2-Pairwise-K16-B64 & DFlash2-Unary-K16-B64 & 0.502 & 0.500 & 4.09 & 8.231 & 8.261 & 0.149 & 22.122 & 22.142 & 13.89 \\
HumanEval & DFlash2-Pairwise-K16-B16 & Raw-DFlash & -1.495 & -1.512 & -8.86 & 8.235 & 4.598 & 1.001 & 22.094 & 21.658 & -84.16 \\
HumanEval & DFlash2-Pairwise-K16-B16 & DFlash-Original-DDTree-B16 & -2.394 & -2.411 & -14.17 & 8.235 & 4.594 & 0.499 & 22.094 & 20.710 & -110.43 \\
HumanEval & DFlash2-Pairwise-K16-B32 & Raw-DFlash & -1.115 & -1.132 & -5.24 & 8.246 & 4.598 & 1.325 & 22.271 & 21.658 & -75.49 \\
HumanEval & DFlash2-Pairwise-K16-B32 & DFlash-Original-DDTree-B32 & -2.711 & -2.728 & -13.48 & 8.246 & 4.590 & 0.782 & 22.271 & 20.827 & -123.26 \\
HumanEval & DFlash2-Pairwise-K16-B64 & Raw-DFlash & -0.831 & -0.850 & -2.99 & 8.231 & 4.598 & 1.950 & 22.122 & 21.658 & -69.53 \\
HumanEval & DFlash2-Pairwise-K16-B64 & DFlash-Original-DDTree-B64 & -2.813 & -2.832 & -12.90 & 8.231 & 4.585 & 1.334 & 22.122 & 20.664 & -131.61 \\
MATH500 & DFlash2-Pairwise-K16-B16 & DFlash2-Original-DDTree-B16 & 0.530 & 0.526 & 4.35 & 7.858 & 7.874 & 0.126 & 21.300 & 21.304 & 15.63 \\
MATH500 & DFlash2-Pairwise-K16-B7 & DFlash2-Original-DDTree-B7 & 0.276 & 0.277 & 2.89 & 7.845 & 7.869 & 0.096 & 21.055 & 21.124 & 8.88 \\
MATH500 & DFlash2-Pairwise-K16-B32 & DFlash2-Original-DDTree-B32 & 0.611 & 0.608 & 4.41 & 7.856 & 7.868 & -0.151 & 21.386 & 21.387 & 19.30 \\
MATH500 & DFlash2-Pairwise-K16-B64 & DFlash2-Original-DDTree-B64 & 0.636 & 0.633 & 4.32 & 7.840 & 7.866 & -1.841 & 21.227 & 21.246 & 28.28 \\
MATH500 & DFlash2-Pairwise-K16-B7 & Raw-DFlash2 & 0.285 & 0.284 & 3.29 & 7.845 & 7.110 & 0.774 & 21.055 & 21.956 & 2.75 \\
MATH500 & DFlash2-Pairwise-K16-B32 & DFlash2-Unary-K16-B32 & 0.616 & 0.613 & 4.48 & 7.856 & 7.865 & 0.137 & 21.386 & 21.359 & 17.73 \\
MATH500 & DFlash2-Pairwise-K16-B64 & DFlash2-Unary-K16-B64 & 0.637 & 0.634 & 4.32 & 7.840 & 7.871 & 0.147 & 21.227 & 21.257 & 18.51 \\
MATH500 & DFlash2-Pairwise-K16-B16 & Raw-DFlash & -1.437 & -1.458 & -5.62 & 7.858 & 4.493 & 0.948 & 21.300 & 21.222 & -84.06 \\
MATH500 & DFlash2-Pairwise-K16-B16 & DFlash-Original-DDTree-B16 & -2.077 & -2.098 & -9.37 & 7.858 & 4.504 & 0.461 & 21.300 & 20.363 & -101.17 \\
MATH500 & DFlash2-Pairwise-K16-B32 & Raw-DFlash & -1.107 & -1.129 & -3.38 & 7.856 & 4.493 & 1.246 & 21.386 & 21.222 & -76.35 \\
MATH500 & DFlash2-Pairwise-K16-B32 & DFlash-Original-DDTree-B32 & -2.427 & -2.450 & -9.91 & 7.856 & 4.495 & 0.721 & 21.386 & 20.401 & -115.55 \\
MATH500 & DFlash2-Pairwise-K16-B64 & Raw-DFlash & -0.885 & -0.908 & -1.84 & 7.840 & 4.493 & 1.833 & 21.227 & 21.222 & -72.09 \\
MATH500 & DFlash2-Pairwise-K16-B64 & DFlash-Original-DDTree-B64 & -2.721 & -2.743 & -10.02 & 7.840 & 4.485 & 1.239 & 21.227 & 20.209 & -130.37 \\
MT-Bench-controlled & DFlash2-Pairwise-K16-B16 & DFlash2-Original-DDTree-B16 & 0.245 & 0.241 & 6.71 & 7.688 & 7.711 & 0.124 & 21.107 & 21.104 & 7.07 \\
MT-Bench-controlled & DFlash2-Pairwise-K16-B7 & DFlash2-Original-DDTree-B7 & 0.161 & 0.157 & 4.84 & 7.699 & 7.701 & 0.099 & 20.885 & 20.901 & 4.65 \\
MT-Bench-controlled & DFlash2-Pairwise-K16-B32 & DFlash2-Original-DDTree-B32 & 0.320 & 0.320 & 7.64 & 7.680 & 7.697 & -0.178 & 21.199 & 21.208 & 10.56 \\
MT-Bench-controlled & DFlash2-Pairwise-K16-B64 & DFlash2-Original-DDTree-B64 & 0.339 & 0.335 & 7.38 & 7.670 & 7.706 & -2.127 & 21.114 & 21.159 & 18.13 \\
MT-Bench-controlled & DFlash2-Pairwise-K16-B7 & Raw-DFlash2 & 0.351 & 0.346 & 15.60 & 7.699 & 6.959 & 0.775 & 20.885 & 21.739 & 5.88 \\
MT-Bench-controlled & DFlash2-Pairwise-K16-B32 & DFlash2-Unary-K16-B32 & 0.307 & 0.305 & 7.41 & 7.680 & 7.682 & 0.138 & 21.199 & 21.175 & 8.70 \\
MT-Bench-controlled & DFlash2-Pairwise-K16-B64 & DFlash2-Unary-K16-B64 & 0.340 & 0.337 & 7.31 & 7.670 & 7.680 & 0.334 & 21.114 & 21.092 & 8.85 \\
MT-Bench-controlled & DFlash2-Pairwise-K16-B16 & Raw-DFlash & 0.088 & 0.066 & 21.16 & 7.688 & 4.448 & 0.959 & 21.107 & 20.975 & -21.65 \\
MT-Bench-controlled & DFlash2-Pairwise-K16-B16 & DFlash-Original-DDTree-B16 & -0.701 & -0.723 & -5.27 & 7.688 & 4.434 & 0.472 & 21.107 & 20.082 & -46.93 \\
MT-Bench-controlled & DFlash2-Pairwise-K16-B32 & Raw-DFlash & 0.406 & 0.384 & 29.18 & 7.680 & 4.448 & 1.275 & 21.199 & 20.975 & -13.67 \\
MT-Bench-controlled & DFlash2-Pairwise-K16-B32 & DFlash-Original-DDTree-B32 & -0.802 & -0.824 & -4.52 & 7.680 & 4.425 & 0.751 & 21.199 & 20.161 & -52.56 \\
MT-Bench-controlled & DFlash2-Pairwise-K16-B64 & Raw-DFlash & 0.652 & 0.628 & 34.64 & 7.670 & 4.448 & 2.087 & 21.114 & 20.975 & -9.46 \\
MT-Bench-controlled & DFlash2-Pairwise-K16-B64 & DFlash-Original-DDTree-B64 & -0.912 & -0.936 & -4.54 & 7.670 & 4.404 & 1.492 & 21.114 & 20.005 & -61.24 \\
MT-Bench-native & DFlash2-Pairwise-K16-B16 & DFlash2-Original-DDTree-B16 & 0.250 & 0.247 & 6.93 & 7.631 & 7.645 & 0.121 & 20.925 & 20.931 & 7.20 \\
MT-Bench-native & DFlash2-Pairwise-K16-B7 & DFlash2-Original-DDTree-B7 & 0.152 & 0.149 & 4.55 & 7.637 & 7.644 & 0.092 & 20.706 & 20.718 & 4.51 \\
MT-Bench-native & DFlash2-Pairwise-K16-B32 & DFlash2-Original-DDTree-B32 & 0.331 & 0.330 & 7.85 & 7.622 & 7.631 & -0.160 & 21.069 & 21.037 & 10.68 \\
MT-Bench-native & DFlash2-Pairwise-K16-B64 & DFlash2-Original-DDTree-B64 & 0.337 & 0.333 & 6.55 & 7.601 & 7.634 & -2.489 & 20.943 & 20.995 & 18.70 \\
MT-Bench-native & DFlash2-Pairwise-K16-B7 & Raw-DFlash2 & 0.345 & 0.340 & 15.64 & 7.637 & 6.905 & 0.758 & 20.706 & 21.536 & 5.68 \\
MT-Bench-native & DFlash2-Pairwise-K16-B32 & DFlash2-Unary-K16-B32 & 0.327 & 0.324 & 8.21 & 7.622 & 7.625 & 0.134 & 21.069 & 21.028 & 9.31 \\
MT-Bench-native & DFlash2-Pairwise-K16-B64 & DFlash2-Unary-K16-B64 & 0.357 & 0.353 & 6.88 & 7.601 & 7.613 & 0.133 & 20.943 & 20.958 & 10.38 \\
MT-Bench-native & DFlash2-Pairwise-K16-B16 & Raw-DFlash & 0.092 & 0.069 & 20.08 & 7.631 & 4.551 & 0.935 & 20.925 & 21.729 & -15.32 \\
MT-Bench-native & DFlash2-Pairwise-K16-B16 & DFlash-Original-DDTree-B16 & -0.707 & -0.729 & -5.67 & 7.631 & 4.529 & 0.441 & 20.925 & 20.801 & -40.29 \\
MT-Bench-native & DFlash2-Pairwise-K16-B32 & Raw-DFlash & 0.427 & 0.405 & 28.62 & 7.622 & 4.551 & 1.243 & 21.069 & 21.729 & -6.95 \\
MT-Bench-native & DFlash2-Pairwise-K16-B32 & DFlash-Original-DDTree-B32 & -0.789 & -0.811 & -3.96 & 7.622 & 4.514 & 0.711 & 21.069 & 20.866 & -44.99 \\
MT-Bench-native & DFlash2-Pairwise-K16-B64 & Raw-DFlash & 0.660 & 0.636 & 33.40 & 7.601 & 4.551 & 1.846 & 20.943 & 21.729 & -2.04 \\
MT-Bench-native & DFlash2-Pairwise-K16-B64 & DFlash-Original-DDTree-B64 & -0.940 & -0.963 & -5.31 & 7.601 & 4.503 & 1.243 & 20.943 & 20.742 & -53.30 \\
\end{longtable}


\begin{longtable}{lllrrrr}
\caption{All directly scored task-quality results.}\label{tab:all-quality}\\
\toprule
Task & Family & Method & B & Evaluated & Passed & Score \% \\
\midrule
\endfirsthead
\toprule
Task & Family & Method & B & Evaluated & Passed & Score \% \\
\midrule
\endhead
\midrule
\multicolumn{7}{r}{Continued on next page}\\
\endfoot
\bottomrule
\endlastfoot
gsm8k & dflash2 & target-only & -- & 128 & 65 & 50.78 \\
gsm8k & dflash2 & dflash2 & -- & 128 & 61 & 47.66 \\
gsm8k & dflash2 & dflash2-original-ddtree & 16 & 128 & 64 & 50.00 \\
gsm8k & dflash2 & dflash2-original-ddtree & 32 & 128 & 63 & 49.22 \\
gsm8k & dflash2 & dflash2-original-ddtree & 64 & 128 & 62 & 48.44 \\
gsm8k & dflash2 & dflash2-original-ddtree & 7 & 128 & 62 & 48.44 \\
gsm8k & dflash2 & pairwise & 16 & 128 & 62 & 48.44 \\
gsm8k & dflash2 & pairwise & 32 & 128 & 60 & 46.88 \\
gsm8k & dflash2 & pairwise & 64 & 128 & 65 & 50.78 \\
gsm8k & dflash2 & pairwise & 7 & 128 & 64 & 50.00 \\
gsm8k & dflash2 & unary & 32 & 128 & 62 & 48.44 \\
gsm8k & dflash2 & unary & 64 & 128 & 62 & 48.44 \\
gsm8k & original & target-only & -- & 128 & 65 & 50.78 \\
gsm8k & original & dflash-original-ddtree & 16 & 128 & 63 & 49.22 \\
gsm8k & original & dflash-original-ddtree & 32 & 128 & 60 & 46.88 \\
gsm8k & original & dflash-original-ddtree & 64 & 128 & 61 & 47.66 \\
gsm8k & original & dflash & -- & 128 & 61 & 47.66 \\
math500 & dflash2 & target-only & -- & 128 & 18 & 14.06 \\
math500 & dflash2 & dflash2 & -- & 128 & 17 & 13.28 \\
math500 & dflash2 & dflash2-original-ddtree & 16 & 128 & 18 & 14.06 \\
math500 & dflash2 & dflash2-original-ddtree & 32 & 128 & 17 & 13.28 \\
math500 & dflash2 & dflash2-original-ddtree & 64 & 128 & 20 & 15.62 \\
math500 & dflash2 & dflash2-original-ddtree & 7 & 128 & 19 & 14.84 \\
math500 & dflash2 & pairwise & 16 & 128 & 18 & 14.06 \\
math500 & dflash2 & pairwise & 32 & 128 & 17 & 13.28 \\
math500 & dflash2 & pairwise & 64 & 128 & 18 & 14.06 \\
math500 & dflash2 & pairwise & 7 & 128 & 16 & 12.50 \\
math500 & dflash2 & unary & 32 & 128 & 17 & 13.28 \\
math500 & dflash2 & unary & 64 & 128 & 20 & 15.62 \\
math500 & original & target-only & -- & 128 & 18 & 14.06 \\
math500 & original & dflash-original-ddtree & 16 & 128 & 17 & 13.28 \\
math500 & original & dflash-original-ddtree & 32 & 128 & 18 & 14.06 \\
math500 & original & dflash-original-ddtree & 64 & 128 & 18 & 14.06 \\
math500 & original & dflash & -- & 128 & 17 & 13.28 \\
humaneval & dflash2 & target-only & -- & 164 & 91 & 55.49 \\
humaneval & dflash2 & dflash2 & -- & 164 & 89 & 54.27 \\
humaneval & dflash2 & dflash2-original-ddtree & 16 & 164 & 88 & 53.66 \\
humaneval & dflash2 & dflash2-original-ddtree & 32 & 164 & 86 & 52.44 \\
humaneval & dflash2 & dflash2-original-ddtree & 64 & 164 & 89 & 54.27 \\
humaneval & dflash2 & dflash2-original-ddtree & 7 & 164 & 91 & 55.49 \\
humaneval & dflash2 & pairwise & 16 & 164 & 89 & 54.27 \\
humaneval & dflash2 & pairwise & 32 & 164 & 90 & 54.88 \\
humaneval & dflash2 & pairwise & 64 & 164 & 91 & 55.49 \\
humaneval & dflash2 & pairwise & 7 & 164 & 88 & 53.66 \\
humaneval & dflash2 & unary & 32 & 164 & 87 & 53.05 \\
humaneval & dflash2 & unary & 64 & 164 & 89 & 54.27 \\
humaneval & original & target-only & -- & 164 & 91 & 55.49 \\
humaneval & original & dflash-original-ddtree & 16 & 164 & 89 & 54.27 \\
humaneval & original & dflash-original-ddtree & 32 & 164 & 88 & 53.66 \\
humaneval & original & dflash-original-ddtree & 64 & 164 & 91 & 55.49 \\
humaneval & original & dflash & -- & 164 & 90 & 54.88 \\
\end{longtable}

\end{landscape}

\end{document}
